% Содержимое подраздела 2.1 - Гетероэнкодеры и вариационные автоэнкодеры: оптимизация в латентном пространстве

\section{Гетероэнкодеры и вариационные автоэнкодеры: оптимизация в латентном пространстве}

Подходы, основанные на автоэнкодерах, преследуют цель создать непрерывное, низкоразмерное и химически осмысленное <<латентное пространство>>, в котором молекулярные структуры представлены в виде векторов. Оптимизация свойств в этом случае сводится к навигации по латентному пространству для поиска векторов, которые после декодирования преобразуются в молекулы с желаемыми характеристиками.

\rgnsubsection{Теоретические основы: концепция гетероэнкодера}

Стандартные автоэнкодеры, обученные на строковом представлении молекул (SMILES), сталкиваются с фундаментальной <<проблемой репрезентации>>. Поскольку одна и та же молекулярная структура может быть описана множеством различных, но химически эквивалентных SMILES-строк, стандартный автоэнкодер обучается сжимать и восстанавливать не инвариантную структуру молекулы, а особенности конкретной текстовой строки. Это приводит к тому, что разные SMILES-представления одной и той же молекулы отображаются в удаленные точки латентного пространства, делая его <<зашумленным>>, неструктурированным и непригодным для эффективной навигации \cite{bjerrum2018improving}.

Для решения этой проблемы была предложена концепция химического \textbf{гетероэнкодера}. В отличие от стандартного автоэнкодера, который обучается на задаче реконструкции (вход идентичен выходу), гетероэнкодер решает задачу <<перевода>> между различными, но семантически эквивалентными представлениями одной и той же сущности. В контексте химии это означает, что модель может обучаться, например, переводить каноническую SMILES-строку в одну из множества случайных (enumerated) SMILES-строк. Такая постановочная задача заставляет модель игнорировать синтаксические особенности входной строки и изучать инвариантное, химически релевантное представление молекулы.

Эффективность данного подхода была продемонстрирована количественно. Стандартный автоэнкодер (`Can2Can`) показал сильную корреляцию расстояний в латентном пространстве со схожестью текстовых строк (коэффициент детерминации $R^2 = 0.58$) и слабую -- со схожестью молекулярных отпечатков ($R^2 = 0.24$). В то же время гетероэнкодер `Can2Enum` (канонический SMILES на входе, случайный на выходе) достиг сбалансированного и химически релевантного пространства, где корреляция с отпечатками достигла $R^2 = 0.58$, а со строками -- $R^2 = 0.55$. Латентные векторы, полученные с помощью гетероэнкодеров, также показали себя значительно более качественными признаками для решения задач QSAR, превзойдя как стандартный автоэнкодер, так и классические отпечатки ECFP4 \cite{bjerrum2018improving}.

\rgnsubsection{Пример реализации: гетероэнкодер для дизайна ингибиторов Bcr-Abl}

Практическое применение и дальнейшее развитие концепции гетероэнкодера было продемонстрировано в задаче \textit{de novo} дизайна потенциальных ингибиторов Bcr-Abl тирозинкиназы --- ключевой мишени в терапии хронического миелоидного лейкоза \cite{карпенко2023генеративная}. В данной работе была разработана и апробирована сложная архитектура условного гетероэнкодера, специально адаптированная для интеграции данных различной природы и целенаправленной генерации молекул с высоким сродством к белковой мишени.

Архитектура модели, детально описанная в работе, представляет собой реализацию мультимодального подхода:
\begin{itemize}
    \item \textbf{Блок кодирования} включает три параллельных энкодера для обработки разнородных входных данных: два LSTM-энкодера для стандартной и канонической SMILES-строк и один полносвязный энкодер для вектора физико-химических дескрипторов. Такой подход позволяет создать максимально полное представление исходной молекулы.

    \item \textbf{Латентное пространство} формируется путем конкатенации выходных векторов всех трех энкодеров. Ключевой особенностью, превращающей модель в условную, является добавление к этому вектору одного нейрона, отвечающего за целевое значение энергии связывания. Это позволяет в дальнейшем управлять процессом генерации, "заказывая" у модели молекулы с желаемым свойством.

    \item \textbf{Блок декодирования} состоит из двух параллельных LSTM-декодеров, которые обучаются восстанавливать из единого условного латентного вектора как стандартную, так и каноническую SMILES-строки, реализуя принцип multi-task learning.
\end{itemize}

Модель была обучена на наборе из 108 410 молекул, содержащих скаффолд 2-ариламинопиримидина. В результате последующей генерации и компьютерной валидации методами молекулярного докинга были отобраны четыре соединения-лидера. Эти соединения продемонстрировали предсказанную аффинность, сопоставимую или превосходящую референсный препарат понатиниб. Например, соединение-лидер I показало энергию связывания с нативной формой Bcr-Abl $-13.8$ ккал/моль, что значительно лучше показателя понатиниба ($-12.0$ ккал/моль) \cite{карпенко2023генеративная}. 

Данный кейс наглядно демонстрирует, как архитектура гетероэнкодера может быть успешно адаптирована для решения конкретной задачи в медицинской химии, реализуя полный вычислительный пайплайн от генерации до валидации кандидатов. Ниже, в следующем подразделе, будет представлен детальный анализ практической реализации этой модели на основе ее открытого исходного кода.

\rgnsubsection{Практический анализ реализации: обзор кода и архитектуры}

Для глубокого понимания теоретической модели, описанной в предыдущем разделе, был проведен статический анализ ее эталонной реализации, представленной в открытом репозитории \texttt{Smiles-HeteroEncoder0-CML}. Данный анализ позволяет детально рассмотреть практические аспекты построения и обучения модели, а также реконструировать пайплайн обработки данных, объясняя не только последовательность шагов, но и их методологическое обоснование.

\rgnsubsection{Технологический стек.}

Проект построен на основе стандартного для задач глубокого обучения и хемоинформатики набора библиотек. Роль каждой из них четко определена:
\begin{itemize}
    \item \textbf{\texttt{tensorflow} (версия 2.4.1):} Выступает в качестве основного фреймворка для построения и обучения нейронной сети. Он предоставляет все необходимые высокоуровневые компоненты через Keras API, включая рекуррентные слои (\texttt{LSTM}), полносвязные слои (\texttt{Dense}), а также механизмы для компиляции модели и управления процессом обучения.
    \item \textbf{\texttt{rdkit-pypi}:} Является краеугольным камнем для решения хемоинформатических задач. Библиотека играет центральную роль в пайплайне обработки данных, предоставляя инструменты для парсинга SMILES-строк в молекулярные графы (\texttt{Chem.MolFromSmiles}), выполнения поиска по подструктуре (\texttt{HasSubstructMatch}) для фильтрации по фармакофору и вычисления молекулярных дескрипторов.
    \item \textbf{\texttt{scikit-learn}:} Используется для выполнения классических задач машинного обучения, которые предшествуют и поддерживают основной deep learning пайплайн. В частности, с помощью этой библиотеки решаются задачи разделения выборки на обучающую и тестовую, а также стандартизация числовых признаков (дескрипторов и энергии связывания) с помощью \texttt{StandardScaler}, что является обязательным шагом для стабилизации обучения нейронной сети.
\end{itemize}

\rgnsubsection{Пайплайн обработки данных.}

Пайплайн обработки данных, реконструированный на основе анализа кода, включает несколько последовательных этапов, каждый из которых имеет четкое назначение.
\begin{itemize}
    \item \textbf{Первичная фильтрация по длине.} Исходные данные, представленные в виде SMILES-строк и связанных с ними свойств, проходят фильтрацию, в ходе которой отбираются только те молекулы, длина SMILES которых находится в диапазоне от 35 до 75 символов. Цель этого шага -- унификация входных последовательностей. Это позволяет уменьшить объем необходимого паддинга (дополнения последовательностей до одинаковой длины) и стабилизировать обучение рекуррентных слоев, которые чувствительны к большой вариативности длин входов.
    \item \textbf{Фокусировка на целевом химическом пространстве.} Далее, с помощью \texttt{rdkit}, каждая молекула проверяется на наличие ключевой фармакофорной группы (2-арилпиримидина). Данная операция позволяет сфокусировать модель на генерации молекул, которые уже содержат структурный мотив, известный своей важностью для целевой биологической активности (ингибирование Bcr-Abl тирозинкиназы).
    \item \textbf{Векторизация SMILES-строк.} Финальным шагом подготовки является преобразование символьных SMILES-строк в числовой формат, понятный нейронной сети. Для этого сначала формируется словарь уникальных символов (токенов), встречающихся в датасете. Этот словарь дополняется специальными токенами: \texttt{!} для обозначения начала последовательности и \texttt{E} -- для ее конца. Затем каждая SMILES-строка преобразуется в числовую матрицу методом One-Hot кодирования. При этом методе каждый символ представляется в виде разреженного бинарного вектора, где все элементы равны нулю, кроме одного, соответствующего индексу символа в словаре. В результате, для батча молекул формируется трехмерный тензор размерностью \texttt{(batch\_size, max\_length, vocab\_size)}, полностью готовый для подачи на вход рекуррентной нейронной сети.
\end{itemize}

\rgnsubsection{Архитектура модели.}

Архитектура модели, реализованная с помощью Keras Functional API, в точности соответствует концепции условного гетероэнкодера. Она состоит из трех параллельных энкодеров, латентного пространства и двух параллельных декодеров.

\begin{enumerate}
    \item \textbf{LSTM-энкодеры для SMILES.} Два идентичных по архитектуре энкодера обрабатывают два различных, но химически эквивалентных представления молекулы: исходную и каноническую SMILES-строки. Каждый энкодер состоит из двух последовательных \texttt{LSTM}-слоев. Их задача -- итеративно прочитать последовательность символов и закодировать всю структурную информацию в векторы финальных скрытых состояний (\texttt{hidden state} и \texttt{cell state}). Использование двух энкодеров для разных представлений является формой ансамблирования, что делает латентное представление более робастным и инвариантным к конкретной форме записи SMILES.

    \item \textbf{Dense-энкодер для дескрипторов.} Полносвязная нейронная сеть обрабатывает вектор числовых молекулярных дескрипторов (например, топологических индексов или фингерпринтов). Архитектура включает несколько \texttt{Dense}-слоев с функцией активации \texttt{ReLU} и слой \texttt{BatchNormalization}. Включение \texttt{BatchNormalization} имеет критическое значение для стабилизации обучения, так как он нормализует активации между слоями, предотвращая проблему затухания и взрыва градиентов. Этот энкодер извлекает компактное признаковое представление из неструктурных данных о молекуле.

    \item \textbf{Формирование условного латентного пространства.} Выходные векторы всех трех энкодеров конкатенируются в единый вектор. Этот объединенный вектор пропускается через полносвязный слой-«бутылочное горлышко» (\texttt{bottleneck}). Роль этого слоя заключается в сжатии всей разнородной информации в единое, компактное и информационно насыщенное латентное представление. Для реализации механизма условной генерации к этому латентному вектору добавляется нормализованное значение энергии связывания, подаваемое на отдельный входной нейрон (\texttt{input\_energy}). Этот шаг превращает модель в условную: он позволяет в дальнейшем управлять свойствами генерируемых молекул, задавая желаемое значение энергии.

    \item \textbf{LSTM-декодеры.} Полученный условный латентный вектор преобразуется с помощью \texttt{Dense}-слоев в начальные скрытые состояния для двух параллельных \texttt{LSTM}-декодеров. Каждый декодер решает задачу авторегрессионной генерации последовательности токенов, восстанавливая одно из представлений SMILES (исходное или каноническое). На выходе каждого декодера стоит \texttt{Dense}-слой с функцией активации \texttt{softmax}, которая преобразует выходные логиты в вероятностное распределение по всему словарю символов. Это позволяет на каждом шаге генерации предсказывать наиболее вероятный следующий символ.
\end{enumerate}

\rgnsubsection{Процесс обучения и генерации.}

Процесс обучения модели реализован с использованием оптимизатора \texttt{Adam} и функции потерь \\ \texttt{categorical\_crossentropy}, которая применяется к каждому токену генерируемой последовательности. Важной особенностью обучения декодеров является применение стратегии \textbf{«Teacher Forcing»}. При этой стратегии на вход декодера на каждом временном шаге подается не его собственное предсказание с предыдущего шага, а истинный токен из целевой последовательности. Это значительно стабилизирует и ускоряет обучение, так как предотвращает накопление ошибок на ранних этапах. Для повышения эффективности обучения используется callback \texttt{ReduceLROnPlateau}. Его назначение -- динамическая подстройка скорости обучения. Он отслеживает функцию потерь на валидационной выборке и автоматически снижает скорость обучения, если прогресс останавливается, что позволяет модели более точно «спуститься» в локальный минимум.

Генерация новых молекул (инференс) представляет собой итеративный авторегрессионный процесс.
\begin{enumerate}
    \item \textbf{Инициализация:} Декодер инициализируется латентным вектором, который может быть получен из молекулы-прототипа или сэмплирован из латентного пространства, а также целевым значением энергии.
    \item \textbf{Старт:} На первом шаге на вход декодера подается стартовый токен (\texttt{!}).
    \item \textbf{Авторегрессионный цикл:} На каждом последующем шаге модель предсказывает следующий токен на основе своего текущего состояния и входного токена с предыдущего шага. Предсказанный токен добавляется к генерируемой строке и одновременно используется в качестве входа для следующего шага. Этот цикл «предсказал-добавил-подал на вход» продолжается итеративно.
    \item \textbf{Завершение:} Процесс останавливается, когда модель генерирует стоп-токен (\texttt{E}) или достигается максимальная заданная длина последовательности.
\end{enumerate}

\rgnsubsection{Развитие VAE-подходов: решение проблем и повышение эффективности}

Вариационные автоэнкодеры (VAE), являясь вероятностной версией автоэнкодеров, стали основным инструментом для генерации молекул. Однако их базовые реализации сталкивались с рядом проблем, что стимулировало разработку более совершенных архитектур.

Одной из главных проблем является \textbf{коллапс апостериорного распределения} (\textit{posterior collapse}), при котором декодер обучается игнорировать информацию из латентного пространства и генерирует молекулы, основываясь только на запомненных данных. Для решения этой проблемы была предложена модель PCF-VAE \cite{bhadwal2025pcf}. Ключевой инновацией стала репараметризация функции потерь ELBO путем введения динамического весового коэффициента $\alpha$ для члена KL-дивергенции. Этот коэффициент постепенно увеличивается в ходе обучения (annealing), позволяя модели на ранних этапах сосредоточиться на качественной реконструкции, а затем -- на регуляризации латентного пространства. В сочетании с использованием упрощенного синтаксиса GenSMILES это позволило значительно повысить новизну генерируемых молекул (93.77\% у PCF-VAE против 69.49\% у стандартного VAE на бенчмарке MOSES).

Следующим шагом стало повышение эффективности \textbf{условной генерации}. Модель $\beta$-CVAE \cite{richards2022conditional} представила элегантный метод явного кодирования свойств в латентное пространство (Explicit Latent Conditioning, ELC). Вместо стандартного априорного распределения $\mathcal{N}(0, I)$, в функции потерь используется условное распределение $\mathcal{N}(\hat{c}, I)$, где $\hat{c}$ -- нормализованное значение целевого свойства. Это <<встраивает>> информацию о свойстве непосредственно в структуру латентного пространства. В сочетании с ПИ-регулятором для коэффициента $\beta$, поддерживающим целевой уровень взаимной информации, данный подход позволил достичь прорывных результатов. В задаче безусловной оптимизации пенализированного LogP (pLogP) модель достигла значения 104.29, установив новый state-of-the-art с огромным отрывом от предыдущего лучшего результата (38.57).

Наконец, для повышения \textbf{робастности и качества генерации вне распределения} (Out-of-Distribution, OOD) была предложена архитектура Multi-Decoder VAE (MD-VAE) \cite{kwon2023string}. Она состоит из одного общего энкодера и ансамбля из $K$ декодеров. Для поощрения разнообразия каждый декодер обучается на своем собственном латентном векторе $z_k$, сэмплированном из одного апостериорного распределения. Ключевой особенностью является коллаборативная функция потерь, которая вычисляется на основе усредненного предсказания всего ансамбля. Этот подход показал систематическое превосходство над базовыми моделями, особенно в OOD-сценарии, где средняя эффективность генерации у MD-VAE составила 47.7\%, что является относительным улучшением на 31.4\% по сравнению с сильной базовой моделью ControlVAE (36.3\%).
